% Galeria rysunków do przeglądu: każdy rysunek w skali, w jakiej trafi do pracy (szerokość kolumny 160 mm),
% z proponowanym podpisem i miejscem w pracy. Kompilacja: latexmk -pdf -outdir=build galeria.tex
\documentclass[12pt,a4paper]{article}
\usepackage[T1]{fontenc}
\usepackage[utf8]{inputenc}
\usepackage[polish]{babel}
\usepackage{times}
\usepackage{graphicx}
\usepackage[lmargin=30mm, rmargin=20mm, tmargin=25mm, bmargin=25mm]{geometry}
\usepackage[font=small, labelfont=bf, labelsep=period]{caption}
\captionsetup[figure]{name=Rys.}
\newcommand{\rys}[3]{%
  \begin{figure}[p]
    \centering
    \includegraphics[width=\linewidth, height=0.8\textheight, keepaspectratio]{#1}
    \caption{#2}
    \smallskip\footnotesize\textit{Plik: \texttt{\detokenize{#1}.pdf}; miejsce w pracy: #3}
  \end{figure}\clearpage}

\begin{document}
\section*{Rysunki do pracy: przegląd}
Schematy (draw.io, źródła w \texttt{zrodla/}) i rysunki metod (matplotlib, skrypty w \texttt{skrypty/}).
Wykresy wyników z rozdziału 10 eksportuje się z dashboardu (zakładka „Porównanie”, przycisk aparatu).
\clearpage

\rys{przeplyw}{Pakiety jednego połączenia TCP tworzą jeden przepływ dwukierunkowy}{2.2}
\rys{izolacja_if}{Losowe podziały w Isolation Forest: anomalię izoluje mniej podziałów niż punkt normalny}{3.1}
\rys{uwaga}{Uwaga iloczynowa skalowana i wiele głowic w bloku Transformera}{3.3}
\rys{domeny}{Domena źródłowa i docelowa: przeniesienie bez adaptacji (S1) i z adaptacją (S2)}{3.4}
\rys{potok}{Potok danych i architektura systemu}{5.2}
\rys{topologia}{Topologia sieci laboratoryjnej w Mininecie (ramka przerywana: hosty biorące udział w atakach)}{6.1}
\rys{krzywa_dobowa}{Mnożnik intensywności sesji użytkowników w ciągu doby i godziny startu sesji M0--M6}{6.3}
\rys{os_czasu}{Przepływy w sesji M1 w przedziałach 10~s na tle okien ataków}{6.5}
\rys{realizm_ruchu}{Odstępy między sesjami użytkowników na tle rozkładu wykładniczego i odstępy heartbeatu IoT}{6.6}
\rys{regula_okna}{Reguła etykiet: okno ataku i końcówka do 2~s po oknie}{7.2}
\rys{flowtransformer}{Budowa FlowTransformera w pracy}{3.3 i 8.1}
\rys{laczenie}{Przygotowanie zbioru źródłowego CIC-IDS2017 + GeNIS}{8.2}
\rys{scenariusze}{Scenariusze S0, S1 i S2 (białe tło: zbiór źródłowy, szare: Mininet)}{8.3}
\end{document}
